% \section{결론, 한계 및 향후 과제}

% \subsection{종합 결론}

% 본 연구는 국방 무인체계 스택 (UAV/UGV/위성-BLOS) 을 대상으로, 기 문서화된 공격 기법들을 오직 방어 체계의 검증 및 강화를 위한 수단으로만 활용하는 방어 지향적 (red/blue) 시뮬레이션 시스템을 ArduPilot SITL 환경 상에 구축하고 실측하였다. 본 연구는 실제 기체, 실제 RF 신호 및 실제 표적을 배제하였으며, 제시된 모든 대표 수치는 원시 데이터 캡처 기반의 재현 가능한 근거만을 바탕으로 산출되었다.

% 본 시스템의 학술적 및 기술적 기여는 다음의 세 가지 축으로 요약된다. 
% 첫째, 도구 근거 에이전트 기반 방어 체계의 확립이다. 3개의 도메인별 RED 계획기 (UAV RED, UGV RED, 위성/BLOS RED) 와 1개의 단일 공유 BLUE 커맨더가 의사결정 시점에 실제 local llm cli 를 호출하며, 결정론적 탐지 도구 (D2--D8, 센서 교차 검증 등) 를 분석 도구로써 활용한다. 이를 통해 위협 탐지는 하위 센서가 전담하고, 상관분석, 귀속 및 대응 계획 수립은 LLM이 수행하는 역할 분리 구조를 실제 코드로 구현하였다.

% 둘째, 실측 기반의 예방 효과 검증이다. 단순 탐지에서 나아가, 단일 공유 BLUE 커맨더의 자체적인 의사결정이 SITL을 실질적으로 제어하는 실시간 폐루프 시스템을 구축함으로써, 공격성공률이 1.0에서 0.0으로 감소함을 인과적으로 입증하였다 (동일한 의사결정 조건 하에서 디스패치 비활성화 시 ASR 1.0, 활성화 시 0.0 기록). 주요 실측 사례는 다음과 같다: RC 오버라이드 탐지 시 제동 (BRAKE) 모드 전환 (53.13 $\to$ 1.86 m), UGV 조향 공격 시 대기 (HOLD) 모드 전환 (180° $\to$ 29.5°), GPS 스푸핑 시 EKF 소스 전환 (39.78 $\to$ 2.91 m), 그리고 명령 주입 대응을 위한 MAVLink 2 서명 게이트웨이 개입 (ASR 1.0 $\to$ 0.0).

% 셋째, 객관적 평가 방법론의 적용이다. 평가 루프 내 LLM의 개입을 완전히 배제한 결정론적 오라클 채점기 (\texttt{oracle\_scorer.py}) 를 도입하였다. 그 결과, 도구를 정상적으로 활용한 에이전트 지휘관 (agentic\_commander) 은 1.0의 재현율 (recall) 을 달성한 반면, 도구가 없는 블라인드 판정 모델은 0.5에 그침을 정량적으로 증명하였다. 나아가, 실환경 조율 캠페인에서 UGV RED 에이전트가 ``자율 (AUTO) 모드가 액추에이터를 격리하고 있으므로, 조향 오버라이드 명령이 서보에 도달하기 위해서는 먼저 AUTO에서 HOLD 모드로의 전환을 강제해야 한다''는 실험 측정 기반 Kill chain을 자율적으로 구성 및 실행하였고, 이에 대응하여 단일 공유 BLUE 커맨더가 공격 절차를 성공적으로 탐지함을 확인하였다.

% 결론적으로 본 시스템은 (i) 실제 LLM 호출 기반으로 구동되는 다중 에이전트 중심의 Red/Blue 아키텍처, (ii) 탐지, 차단, 복구의 전 과정이 실증적 근거로 연계된 폐루프 방어 메커니즘, 그리고 (iii) LLM을 배제한 결정론적 채점 기반의 객관적 성능 검증을 단일화된 재현 가능 패키지로 제시한다. 본 논문에서 제시된 모든 정량적 수치는 실측값을 바탕으로 하며, 미측정되었거나 설계 단계에 머무른 항목은 그 상태를 투명하게 명시하는 \textbf{evidence-honesty tagging} 을 도입하여 연구 결과의 무결성을 확보하였다.

% \subsection{근본 한계}

% 본 연구에서 제안하는 시스템의 근본적인 한계점을 투명하게 명시한다. 아래에 기술된 한계 중 상당수는 향후 실증 연구를 통해 극복되어야 할 영역이며, 현재의 시뮬레이션 범위 내에서는 해결하기 어려운 구조적 제약 사항에 해당한다.

% \textbf{ (1) 시뮬레이션 전용 및 실기체 미검증:} 모든 공격 및 방어 시나리오는 ArduPilot SITL 환경에서만 수행되었다. 즉, 실제 항공기, 차량, 위성 링크, 실제 RF 재밍 및 스푸핑 하드웨어, 그리고 물리적 표적은 실험에 사용되지 않았다. 위성/BLOS C2 구조의 경로상 노출 조건은 커널 권한 없이 유저 공간에서 동작하는 중계기 에뮬레이션으로 구현되었으며, 물리 계층의 전파 특성, 안테나, 링크 예산 등은 모델링되지 않았다. 실측된 예방 수치 (예: 제동 후 1.86 m) 는 SITL의 물리 모델에 국한된 결과이므로, 실기체에서의 재현 여부는 향후 HIL 및 실증 시험을 통해서만 확인이 가능하다.

% \textbf{ (2) 소표본 조건 ($n \le 5$):} 실제 local LLM CLI 호출 비용 및 SITL 실행 시간의 제약으로 인해, 대부분의 실시간 폐루프 제어 및 판정 실험은 $n=3$에서 최대 $n=5$ 수준의 소규모 표본을 바탕으로 진행되었다. 따라서 본 결과는 방어 기작의 효과 유무와 그 방향성을 입증하는 인과적 근거로 기능할 뿐, 통계적 유의성이나 신뢰구간을 주장할 수 있는 대규모 데이터셋은 아님을 밝힌다.

% \textbf{ (3) LLM 판정의 잔여 편향성:} 오라클 채점기를 도입하여 탐지 축에 존재하는 `LLM이 채점하는 LLM 편향' 을 상당 부분 제거하였으나, 독립 판정 실험은 여전히 인간 평가자가 아닌 LLM에 의존하는 구조적 한계를 지닌다. 판정 그룹 간에 모델 가중치가 완전히 고정되지 않은 지점이 존재하므로 모델 교란의 가능성이 남아 있다. 실제로 독립 판정 과정에서 초기의 가설이 그대로 재현되지 않음을 정직하게 교정 및 보고하였다. 즉, 도구가 없는 도구 미사용 방어군 역시 시나리오를 최고 수준으로 해결하였으며, 이에 따라 도구의 진정한 가치는 탐지 여부의 결정적 요인이 아닌 판정의 일관성, 결정성, 감사 가능성, 그리고 플레이북 매핑에 있음을 재규정하였다.

% \textbf{ (4) 탐지와 예방의 불일치 및 탐지 전용 ([DETECT-ONLY]) 명시:} 1.0에서 0.0으로의 완벽한 공격성공률 (ASR) 억제 효과는 예방 기작이 라이브로 배선된 특정 공격에서만 성립한다. 그 외의 경우, 탐지 기능은 단지 탐지 소요 시간을 단축할 뿐 ASR 자체를 변화시키지는 못한다. 특히 페일오버를 수행할 건전한 대체 소스가 부재한 자력계, IMU, 기압계 대상의 스푸핑 방어는 근거 정직성 태깅 원칙에 따라 철저히 탐지 전용 (\textbf{[DETECT-ONLY]}) 으로 명기하였다. ADS-B 팬텀 표면 공격 역시 이에 해당한다.

% \textbf{ (5) 탐지 성능의 근본적 하한선 존재:} 내부 적대적 점검 결과 탐지기 성능을 강화하였음에도 불구하고, 약 0.025 m/s 이하의 속도로 서서히 진행되는 저속 GPS 스푸핑 공격은 강화 이후에도 센서의 자연 잡음에 묻혀 식별이 불가하였다. 이처럼 시계열 센서 데이터 기반 탐지에는 물리적 원리에 기인한 탐지 하한선이 존재함을 인정한다.

% \textbf{ (6) 연구 범위의 규율 및 신규 무기화 배제:} 본 제출물은 철저한 방어 목적의 연구이므로, 탐지 회피를 위한 신규 공격 기법을 고안하거나, 회피 알고리즘을 실제 무기 체계 수준으로 최적화하거나, 운용자의 시야를 교란하는 위조 기법 등을 의도적으로 배제하였다. 설계 수준 독트린형 킬체인 및 실측 복합 공격 시나리오들은 기 실측된 원시 공격 프리미티브의 정직한 조합에 불과하며, 새로운 취약점을 무기화하려는 시도가 아님을 명확히 한다.

% \textbf{ (7) 결과 도출에 실패한 실험 항목:} 배터리 페일세이프 스푸핑 방어의 경우, D-BATT 탐지기가 0.665초 만에 오탐 없이 경보를 발화하였으나 실질적인 억제 효과는 측정되지 않았다. 또한 거리계 스푸핑 실험은 시스템 지연 문제로 인해 폐기되었다. 이 두 항목은 본 논문의 어떠한 대표 수치 통계에도 포함되지 않았다.

% \subsection{현실성 심화 및 SITL 내장 동역학·센서 모델을 활용한 효과 수준 현실성 심화}

% 본 연구의 자체 검토 과정에서, ArduPilot SITL이 제공하는 고충실도의 현실적 공격 및 기준선 메커니즘을 충분히 활용하지 못한 6개의 한계 지점을 식별하였다. 이 중 5건에 대해서는 SITL 환경에서 재측정을 수행하여 보완하였으며, 1건 (바람을 인가한 기준선 환경) 은 명확한 음성 결과로 남겨두었다. 재측정 결과들은 단순히 공격의 강도를 높이는 것이 아니라 현상의 현실성을 강화하는 방향으로 도출되었으며, 그 과정에서 드러난 한계점도 함께 기술한다. 이는 본 연구가 상정한 현실성 포락선의 상한을 명확히 규명하는 작업이다.

% \textbf{GNSS 재밍 메커니즘의 고도화:} 기존의 \texttt{gnss\_jam.py}는 \texttt{SIM\_GPS1\_ENABLE=0} 매개변수를 활용한 단일 킬스위치 방식으로, GPS 고정 상태를 6에서 1로 불연속적으로 강하시켰다. 이를 ArduPilot 내장 내비게이션 교란 모델인 \texttt{SIM\_GPS1\_JAM=1} 로 교체한 결과, 위성 수, 측정 속도, 추정 위치가 매 임의 시점마다 요동치며 45초 동안 fix 유형이 92회 전이하는 `불안정하게 흔들리다 완전히 신호를 상실하는' 현실적인 열화 패턴이 재현되었다. 이러한 고충실도 열화 환경하에서도 결정론적 탐지 도구인 D7 은 0.3\,\text{s} 만에 위협을 포착하였으며 (도구 미사용 방어군 기준 무풍 오탐 0건), ArduPilot 자체 EKF 페일세이프에 의한 강제 착륙 (LAND) 명령은 4.88\,\text{s} 에 발동되었다 (\texttt{docs/evidence/enh\_gnss\_jam\_realistic/}).

% \textbf{ADS-B 저속 램프-인 공격의 실측 및 취약점 폐쇄:} 제5장 (T-UAV-8) 에서 자가 진단한 취약점인 ``느린 속도로 접근하는 유령 항적은 기존 규칙 R1을 회피할 가능성이 존재함'' 을 검증하기 위해, 최초 관측 거리 2,497\,\text{m} 에서 12\,\text{m/s} 의 속도로 선형 접근하는 유령 항적 시나리오를 구축하여 실측을 진행하였다. 실험 결과, 해당 항적은 기존의 R1, R2, R3 규칙을 모두 우회하였다 (경보 발화 0/205건). 이러한 방어 공백을 폐쇄하기 위해 지속 접근 규칙인 R4 (sustained-inbound) 를 방어 플러그인에 추가하였으며, 그 결과 근접 경계 밖에서 획득되었으나 지속적으로 순접근하여 경보 포락선에 진입하는 적대적 항적을 105\,\text{s} (상대 거리 1,238.6\,\text{m}) 시점에 정확히 포착하였다. 반면, 단순히 교차 통과하는 정상 항적에 대해서는 오탐 0건을 유지하였다 ( \texttt{docs/evidence/surf\_adsb\_rampin/}, \texttt{src/defenses/adsb\_xcheck.py}).

% \textbf{GPS Seamless Takeover (meaconing) 의 구현:} 균일 램프 주입 방식 대신 Psiaki와 Humphreys~\cite{psiaki_humphreys_2016} 가 제시한 정합 후 분기 메커니즘을 \texttt{gps\_spoof.py}에 반영하였다. 초기 30초 동안은 오프셋을 0으로 설정하여 실제 신호와 정합시키고, 이후 1\,\text{m/s} 의 속도로 40\,\text{m} 까지 서서히 위치를 분기시켰다. 정합 구간 내내 결정론적 탐지 도구 D3 은 침묵하였으나 (오탐 0건, innovation 평균 수치는 0.005로 트립 임계치 0.03의 약 1/6 수준), 분기가 개시되면 D3는 분기 시점으로부터 0.81\,\text{s} 만에 경보를 발화하였다. 이는 기존 균일 램프 주입 시의 TTD ($0.81\,\text{s}$) 와 정확히 일치하는 결과이다. 즉, seamless takeover 기법은 본질적인 탐지 회피가 아니라 경보의 지연 효과만을 가짐이 입증되었다. 위치의 불연속적 도약이 없으므로 절대 경보 시점은 공격자가 설정한 정합 시간 창만큼 유예되지만, 관성 데이터와 innovation 간의 교차 검증 구조에는 분기 거동이 그대로 노출된다. 부수적으로 ArduPilot의 \texttt{SIM\_GPS1\_NOISE} 매개변수가 수평 지터가 아닌 고도 축 전용 오차임을 소스 코드 수준에서 규명하고, 임의의 수평 지터를 인위적으로 합성하는 대신 물리적인 수직 왜곡 효과를 그대로 보고하였다.

% \textbf{IMU 주입 진폭 상한 설정:} 물리적 한계를 무시한 무제한 램프 입력 대신, 음향 공진 기반 IMU 기만 문헌의 실측 상한선을 적용하여 캡을 인가하였다: 자이로스코프 0.10\,\text{rad/s}~\cite{tu2018injected}, 가속도계 0.98\,\text{m/s}$^2$~\cite{trippel2017walnut}. 상한을 초과하는 주입 요청은 소스 코드 수준에서 천장값으로 클램프하도록 설계하였다. 동일한 기동 타이밍에서 이처럼 클램프된 현실적 주입 하의 센서 융합 자세 오차는 평균 55.4$^\circ$ 로 측정된 반면, 무제한 주입 시에는 173.7$^\circ$ 로 나타나 위협도가 약 3.1배 완화됨을 확인하였다. 다만 캡이 적용된 공격 시나리오에서도 기체는 15$^\circ$ 의 임무 실패선을 초과하고 평균 고도가 약 17\,\text{m} 손실되어 여전히 유효한 타격을 입었다. 하위 탐지 도구인 D-IMU는 캡이 적용된 공격에 대해서도 100\% 의 탐지율을 달성하였으나, TTD는 2.11\,\text{s} 를 기록하여 신호 대 잡음비가 뚜렷했던 무제한 주입 시보다 지연되었다. 이는 현실적인 공격일수록 물리적 위협도는 낮아지나 탐지 시점은 다소 지연된다는 트레이드오프 관계를 실증한다. 한편 실험에는 \texttt{SIM\_GYR*\_RND} 및 \texttt{SIM\_ACC*\_RND} 노이즈 모델을 동시 인가하였으며, 실제 파일 기반 주입은 향후 연구 과제로 남겨둔다.

% \textbf{기압 WCF 물리 벡터 모델 적용:} 임의의 고도 오프셋을 강제하는 \texttt{SIM\_BARO\_GLITCH} 대신, 정압구의 기류 교란을 물리적으로 모델링하는 \texttt{WCF\_FWD} 계수와 \texttt{SIM\_BARO\_DELAY} 매개변수를 통해 기만 행위를 묘사하였다. 본 수식 항은 동압에 비례하므로 정지 체공 환경에서는 교란 수치가 0에 수렴하여 (체공 널 대조군 변위 0.31\,\text{m}, D-BARO 침묵), 위협이 은닉된다. 그러나 전진 비행 시 (체감 대기속도 약 8.7\,\text{m/s}) 에는 최대 11.58\,\text{m} 의 고도 추정 오차를 유발한다. 실측된 물리 변위는 해석학적 예측값의 0.96배 수준으로 일치하였으며, 입력 계수 변화에 대해 엄격한 선형성을 보여 물리 모델이 SITL 내부 동역학 구조와 유기적으로 결합되어 있음을 교차 확증하였다. 또한 해당 주입은 GLITCH 매개변수와 부호가 반대로 작용하여 기체의 고도를 강제 상승시키는 특성을 보였다. 정직한 물리적 한계로서 계수 1 수준의 현실적 공격은 임무 실패 기준선인 5\,\text{m} 에 미달하였으며, 10\,\text{m} 이상의 치명적 드리프트를 유발하기 위해서는 극단적인 계수 설정과 고속 비행 조건이 동시에 요구됨이 밝혀졌다. 하위 도구인 D-BARO는 해당 공격을 100\% 의 성공률로 탐지하였다 (TTD 5.96\,\text{s}, 오탐 0건).

% \textbf{바람 환경을 혼합한 기준선의 미해결 상태 고지:} 본 연구의 ``오탐 0건'' 관련 모든 정량적 명제들은 외부 교란이 통제된 무풍 정적 SITL 환경에서 측정된 지표들이다. 이에 따라 탐지 알고리즘의 대지 강건성을 검증하고자 \texttt{SIM\_WIND\_SPD/DIR/TURB} 매개변수를 활용하여 10\,\text{m/s} 의 풍속과 0.3 등급의 난류 환경을 주입하고 오탐율을 재계측하였다. 그러나 시뮬레이터 상에서 바람 매개변수의 인가가 정상 확인되고 고도 프로파일을 상수로 고정하여 재시도하였음에도 불구하고, 위치 유지 중인 기체의 대지 속도, 표류 변위, 동체 기울기는 무풍 상태와 질적으로 동일하였다. 즉, 기저 ArduPilot SITL 멀티로터 물리 모델은 체공 상태에서 주변 정상풍에 의한 공기역학적 항력을 거의 받지 않도록 설계되어 있어, 바람이 기체 프레임에 실질적인 동역학적 교란을 가하지 못하였다. 결과적으로 ``오탐 제로의 성과가 조용한 시뮬레이션 환경에 종속된 아티팩트인가'' 라는 의문은 미해결 영역으로 남았으며, 본 연구진은 상대 대기속도가 충분히 확보되는 전진 비행 영역에서의 재측정을 명확한 향후 과제로 고지한다. 미검증된 강건성을 가설만으로 논문에 계상하지 않는다는 원칙을 고수한다.

% \subsection{심화 관측 및 자체 레드티밍 결과}

% 상기 식별된 현실성 강화 요소들을 AI 에이전트 인프라와 결합하여, 임무 수준에서의 최종 파급 효과와 자체 구축한 결정론적 탐지 도구들의 사각지대를 정밀 측정하였다.

% \begin{itemize}[leftmargin=*]
% %    \item \textbf{GNSS 자율 비행 도메인 심화:} 고충실도 재밍 모델 하에서 기체는 경로를 급격히 이탈하지 않고 오차 포락선 내에서 경합을 유지하였다. 단거리 3개 웨이포인트 미션의 경우, EKF 분산 수치가 페일세이프 임계치를 트리거하기 전 (12.1\,\text{s}) 에 비행 계획을 완주하였으며, 실질적인 피해는 임무 종료 후 강제 발동된 페일세이프 LAND 메커니즘에 의한 비지령 하강과 사후 고도 추정값의 발산 (112\,\text{m}) 으로 제한되었다. 최상위 단일 공유 BLUE 커맨더가 참조하는 D7 도구는 EKF 페일세이프보다 11.65\,\text{s} 선행한 0.47\,\text{s} 시점에 이를 조기 포착하였다.

%     \item \textbf{GNSS 자율 비행 도메인 심화:} 고충실도 재밍 모델 하에서 기체는 경로를 급격히 이탈하지 않고 오차 포락선 내에서 경합을 유지하였다. 단거리 3개 WP 미션의 경우, EKF 분산 수치가 페일세이프 임계치를 트리거하기 전 (12.1\,\text{s}) 에 비행 계획을 완주하였으며, 실질적인 피해는 임무 종료 후 강제 발동된 페일세이프 LAND 메커니즘에 의한 비지령 하강과 사후 고도 추정값의 발산 (112\,\text{m}) 으로 제한되었다. 최상위 단일 공유 BLUE 커맨더가 참조하는 D7 도구는 EKF 페일세이프보다 11.65\,\text{s} 선행한 0.47\,\text{s} 시점에 이를 포착하였다.
    
%     \item \textbf{자체 레드티밍을 통한 맹점 규명 (D7):} 건강성 지표 기반의 D7 탐지 도구는 지연 유발형 공격에 치명적인 사각지대를 가짐이 확인되었다. 수신되는 fix 데이터의 구조적 무결성이 유지되어 D7이 완전히 침묵하는 사이, 기저 EKF 내부 오차는 26.5\,\text{m} 까지 은밀하게 누적되었다. 이처럼 건강한 fix 상태의 내부 부패는 오직 이노베이션 채널을 감시하는 D3 도구에 의해서만 포착되었다. 이를 통해 데이터 건강성 탐지와 물리적 이노베이션 탐지는 구조적으로 직교하며, 상호 대체 불가능한 독립 자원임이 입증되었다.
    
%     \item \textbf{IMU 동적 진동 특성 분석:} 직류 편향을 가해 기체를 한쪽 방향으로 지속적으로 기울이던 기존 근사 모델과 달리, 실제 음향 공진을 모사한 영평균 진동 변위를 주입한 결과 동체는 일방향 기울어짐이 아닌 극심한 고주파 흔들림 거동을 보였다. 이는 실제 기만 공격이 유발하는 동역학적 특성이 단순 직류 오프셋 모델과 질적으로 상이함을 보여준다.
    
%     \item \textbf{기압 자율 비행 (AUTO) 도메인 심화:} \texttt{WCF\_FWD} 모델을 고도 변경이 포함된 AUTO 임무 중에 주입한 결과, 기체는 설정된 목표 고도 대비 11.5\,\text{m} 실제 상승 기동을 수행하였으나 운용자 제어창에 보고되는 피드백 고도는 오직 0.65\,\text{m} 만 변위되어 물리적 위험을 은닉하는 효과를 입증하였다. 측정치와 해석적 예측값의 비율은 0.94로 정합하여 물리 법칙에 부합함을 재확인하였고, 하위 탐지 도구는 공격 개시 후 3.27\,\text{s} 에 위험을 포착하여 단일 공유 BLUE 커맨더에 정보를 환류하였다.
    
% %    \item \textbf{ADS-B R4 규칙의 사각지대 발견:} 추가된 지속 접근 규칙인 R4를 대상으로 자체 레드티밍을 수행한 결과, 정당하게 접근하는 수렴 교통 환경에서 오탐이 유발되는 한계가 관찰되었다. 예를 들어 500\,\text{m} 의 안전 이격 거리를 두고 활주로에 정당하게 정렬 및 착륙 접근하는 실제 기체가 탐지 구역에 진입 시, 단지 순접근율 조건만을 충족한다는 이유로 R4 규칙이 오작동하였다. 이는 알고리즘이 최근접점 및 방위각 변화율 게이트를 누락하여 발생한 아키텍처적 결함으로, 본 연구진은 이에 대한 수정 및 보완을 향후 과제로 적시한다.

%     \item \textbf{ADS-B R4 규칙의 사각지대 발견:} 추가된 지속 접근 규칙인 R4를 대상으로 자체 레드티밍을 수행한 결과, 정당하게 접근하는 수렴 교통 환경에서 오탐이 유발되는 한계가 관찰되었다. 예를 들어 500\,\text{m} 의 안전 이격 거리를 두고 활주로에 정당하게 정렬 및 착륙 접근하는 실제 기체가 탐지 구역에 진입 시, 단지 순접근율 조건만을 충족한다는 이유로 R4 규칙이 오작동하였다. 이는  아키텍처적 결함으로, 본 연구진은 이에 대한 수정 및 보완을 향후 과제로 적시한다.
    
%     \item \textbf{하한 미만 미코닝 (Sub-floor Meaconing) 의 검증 제한:} D3 하드닝 임계치 하한 ($0.024\,\text{m/s}$) 보다 미세하게 설계된 $0.02\,\text{m/s}$ 분기율 조건에서의 물리 실증 실험은 시뮬레이션 런타임 중단 현상으로 인해 최종 완료하지 못하였다. 다만 분기율이 탐지기의 원리적 물리 하한선 미만이므로 D3 도구를 완전히 우회할 것으로 논리적으로 예측되며, 이는 앞서 명시한 탐지 하한 한계 특성과 완벽히 부합한다.
% \end{itemize}

% \subsection{향후 과제}

% \textbf{ (1) 다중 이동체 및 군집 확장}: 현재의 시뮬레이션 환경은 도메인별 단일 이동체만을 대상으로 하고 있다. 향후에는 다수의 UAV 및 UGV가 상호 협조하는 군집 시나리오로 확장하여, 공격 측의 다중 표적 동시 공격 기법과 방어 측의 이동체 간 상호 교차 검증 메커니즘을 구현하고자 한다. 이는 기존에 제안한 도구들의 소스 대 소스 게이팅 사상을 기체 대 기체 차원의 분산 아키텍처로 일반화하는 계기가 될 것이다.

% \textbf{ (2) 실기체 및 HIL 실증}: SITL 환경에서의 실험 결과를 실제 비행 컨트롤러 (HIL) 및 제한된 범위의 실기체 시험으로 이관하여, 실제 물리 동역학 및 적대적 작전환경 조건 하에서 예방 수치의 재현성을 엄밀하게 검증할 필요가 있다. 특히 위성 및 가시권외 (BLOS) 데이터링크 시나리오의 경우, 유저 공간 에뮬레이션 단계를 실제 링크 예산, 신호 지연 및 패킷 손실 모델로 대체하는 것이 최우선 과제이다.

% \textbf{ (3) 독립적 인간 평가}: LLM 판정 메커니즘에 잔존할 수 있는 잔여 편향을 완전히 배제하기 위해, 도메인 전문가에 의한 블라인드 채점 방식과 통제 실험을 도입해야 한다. 또한 전체 실험의 표본 크기를 대폭 확대함으로써, 통계적으로 증명하는 정량적 근거로 격상시킬 필요가 있다.

% \textbf{ (4) 탐지 전용 격차 해소}: 자력계, IMU, 기압계 및 저속 GPS 하한선 문제는 단순한 탐지 임계치의 조정만으로는 근본적인 해결이 불가능하다. 따라서 직교 관측 체계 기반의 물리적 예방 메커니즘을 검증함으로써, 기존의 탐지 전용 항목들을 실질적인 예방 가능 항목으로 전환하는 연구를 추진할 예정이다.

\section{결론 및 향후 과제}

\subsection{결론}

본 연구는 국방 무인체계 스택(UAV, UGV, 위성-BLOS)을 대상으로, 기 문서화된 공격 기법을 방어 체계 검증에 활용하는 방어 지향적(Red/Blue) 시뮬레이션 시스템을 ArduPilot SITL 환경에 구축하고 그 실효성을 실증하였다. 본 시스템의 핵심 학술적, 기술적 기여는 다음 세 가지로 요약된다.

첫째, \textbf{도구 근거 에이전트 기반 방어 체계의 확립}이다. 위협 탐지는 결정론적 하위 센서(D2~D8)가 전담하고, 상황의 상관분석, 귀속 및 대응 계획 수립은 LLM 기반의 단일 공유 BLUE 커맨더가 수행하는 역할 분리 아키텍처를 구현하였다. 

둘째, \textbf{실측 기반의 실시간 폐루프(Closed-loop) 예방 효과 검증}이다. 단순 탐지를 넘어 커맨더의 의사결정이 물리적 제어로 즉각 연동되는 시스템을 통해 공격 성공률(ASR)이 1.0에서 0.0으로 완벽히 억제됨을 입증하였다. 구체적으로 RC 오버라이드 탐지 시 제동(BRAKE) 모드 전환으로 밀림 거리를 53.13m에서 1.86m로 대폭 축소하였으며, GPS 스푸핑 시 EKF 소스 전환, 명령 주입 방어를 위한 MAVLink 2 서명 게이트웨이 개입 등 실효적인 차단 기작을 확인하였다.

셋째, \textbf{객관적 평가 방법론의 도입}이다. 평가 과정에서 LLM의 개입을 완전히 배제한 결정론적 오라클 채점기를 도입하였다. 이를 통해 도구를 정상적으로 활용한 지휘 에이전트(agentic\_commander)가 재현율 1.0을 달성하는 반면, 도구가 없는 판정 모델은 0.5에 그침을 정량적으로 증명하여 도구 결합 아키텍처의 우수성을 입증하였다.

결론적으로 본 연구는 다중 에이전트 중심의 의사결정 체계와 실증적 근거 기반의 폐루프 방어 메커니즘을 융합함으로써, 차세대 자율 무인체계 보안을 위한 결정론적이고 재현 가능한 새로운 방어 방법론을 제시한다.

\subsection{향후 과제}

본 연구의 성과를 바탕으로 다음의 세 가지 발전 방향을 제안한다.

\textbf{(1) 다중 이동체 및 군집(Swarm) 환경으로의 확장}: 단일 이동체를 넘어 다수의 UAV 및 UGV가 협조하는 군집 시나리오를 구축하고, 기체 간 상호 교차 검증 및 분산 방어 아키텍처를 구현한다.

\textbf{(2) 실기체 및 HIL(Hardware-In-the-Loop) 실증}: SITL 환경의 검증 결과를 실제 비행 컨트롤러와 물리적 동역학, 실제 RF 통신(위성 데이터링크 포함) 환경이 반영된 실기체 시험으로 이관하여 방어 체계의 현실적 강건성을 입증한다.

\textbf{(3) 탐지 전용(Detect-only) 영역의 예방 기작 개발}: 자력계, IMU, 기압계 대상의 정밀 스푸핑 등 현재 탐지 수준에 머무르고 있는 위협들에 대해, 독립적인 직교 관측 체계를 융합하여 실질적인 제어권 보전 및 무효화(Mitigation) 메커니즘을 확립한다.